% ============================================================
% ANHANG A — BEGLEITMATERIAL (Video)
% ------------------------------------------------------------
% ZIELDATEI: content/a1_begleitmaterial.tex
% Einbindung in document.tex NACH \appendix:
%   \appendix
%   % ============================================================
% ANHANG A — BEGLEITMATERIAL (Video)
% ------------------------------------------------------------
% ZIELDATEI: content/a1_begleitmaterial.tex
% Einbindung in document.tex NACH \appendix:
%   \appendix
%   % ============================================================
% ANHANG A — BEGLEITMATERIAL (Video)
% ------------------------------------------------------------
% ZIELDATEI: content/a1_begleitmaterial.tex
% Einbindung in document.tex NACH \appendix:
%   \appendix
%   % ============================================================
% ANHANG A — BEGLEITMATERIAL (Video)
% ------------------------------------------------------------
% ZIELDATEI: content/a1_begleitmaterial.tex
% Einbindung in document.tex NACH \appendix:
%   \appendix
%   \include{content/a1_begleitmaterial}
%
% BENOETIGT: fuer den QR-Code das Paket \usepackage{qrcode} im Preamble.
%   (qrcode ist in TeX Live enthalten, kein externes Tool noetig.)
% PLATZHALTER: [VIDEO-URL] durch die stabile, langfristig erreichbare
%   Adresse ersetzen (Uni-Cloud/Repository). Denselben Wert an BEIDEN
%   Stellen eintragen (Text-Link und \qrcode{...}).
% ============================================================

\chapter{Begleitmaterial}
\label{anh:begleitmaterial}

\section{Demonstrationsvideo}
\label{anh:video}

Als Begleitmaterial zu dieser Arbeit steht eine Videoaufzeichnung des Demonstrators zur Verfügung.
Sie zeigt eine frühe Ausbaustufe, in der die Szene über ein Menü in der virtuellen Umgebung manipuliert wird, das vordefinierte Befehle bereitstellt.
Das Video ist unter der folgenden Adresse abrufbar.

\begin{center}
\url{https://git.informatik.fh-nuernberg.de/jobboerse/forschung/bachelorarbeiten/bachelorarbeit-sandro-jaeger/-/blob/fe68f112f06df9da38eff687fd8df8460061186c/BA_VR_Demo_with_menu.mp4}

\vspace{1em}
\qrcode[height=3cm]{https://git.informatik.fh-nuernberg.de/jobboerse/forschung/bachelorarbeiten/bachelorarbeit-sandro-jaeger/-/blob/fe68f112f06df9da38eff687fd8df8460061186c/BA_VR_Demo_with_menu.mp4}

\vspace{0.5em}
{\footnotesize Abbildung: QR-Code zur Videoaufzeichnung.}
\end{center}
%
% BENOETIGT: fuer den QR-Code das Paket \usepackage{qrcode} im Preamble.
%   (qrcode ist in TeX Live enthalten, kein externes Tool noetig.)
% PLATZHALTER: [VIDEO-URL] durch die stabile, langfristig erreichbare
%   Adresse ersetzen (Uni-Cloud/Repository). Denselben Wert an BEIDEN
%   Stellen eintragen (Text-Link und \qrcode{...}).
% ============================================================

\chapter{Begleitmaterial}
\label{anh:begleitmaterial}

\section{Demonstrationsvideo}
\label{anh:video}

Als Begleitmaterial zu dieser Arbeit steht eine Videoaufzeichnung des Demonstrators zur Verfügung.
Sie zeigt eine frühe Ausbaustufe, in der die Szene über ein Menü in der virtuellen Umgebung manipuliert wird, das vordefinierte Befehle bereitstellt.
Das Video ist unter der folgenden Adresse abrufbar.

\begin{center}
\url{https://git.informatik.fh-nuernberg.de/jobboerse/forschung/bachelorarbeiten/bachelorarbeit-sandro-jaeger/-/blob/fe68f112f06df9da38eff687fd8df8460061186c/BA_VR_Demo_with_menu.mp4}

\vspace{1em}
\qrcode[height=3cm]{https://git.informatik.fh-nuernberg.de/jobboerse/forschung/bachelorarbeiten/bachelorarbeit-sandro-jaeger/-/blob/fe68f112f06df9da38eff687fd8df8460061186c/BA_VR_Demo_with_menu.mp4}

\vspace{0.5em}
{\footnotesize Abbildung: QR-Code zur Videoaufzeichnung.}
\end{center}
%
% BENOETIGT: fuer den QR-Code das Paket \usepackage{qrcode} im Preamble.
%   (qrcode ist in TeX Live enthalten, kein externes Tool noetig.)
% PLATZHALTER: [VIDEO-URL] durch die stabile, langfristig erreichbare
%   Adresse ersetzen (Uni-Cloud/Repository). Denselben Wert an BEIDEN
%   Stellen eintragen (Text-Link und \qrcode{...}).
% ============================================================

\chapter{Begleitmaterial}
\label{anh:begleitmaterial}

\section{Demonstrationsvideo}
\label{anh:video}

Als Begleitmaterial zu dieser Arbeit steht eine Videoaufzeichnung des Demonstrators zur Verfügung.
Sie zeigt eine frühe Ausbaustufe, in der die Szene über ein Menü in der virtuellen Umgebung manipuliert wird, das vordefinierte Befehle bereitstellt.
Das Video ist unter der folgenden Adresse abrufbar.

\begin{center}
\url{https://git.informatik.fh-nuernberg.de/jobboerse/forschung/bachelorarbeiten/bachelorarbeit-sandro-jaeger/-/blob/fe68f112f06df9da38eff687fd8df8460061186c/BA_VR_Demo_with_menu.mp4}

\vspace{1em}
\qrcode[height=3cm]{https://git.informatik.fh-nuernberg.de/jobboerse/forschung/bachelorarbeiten/bachelorarbeit-sandro-jaeger/-/blob/fe68f112f06df9da38eff687fd8df8460061186c/BA_VR_Demo_with_menu.mp4}

\vspace{0.5em}
{\footnotesize Abbildung: QR-Code zur Videoaufzeichnung.}
\end{center}
%
% BENOETIGT: fuer den QR-Code das Paket \usepackage{qrcode} im Preamble.
%   (qrcode ist in TeX Live enthalten, kein externes Tool noetig.)
% PLATZHALTER: [VIDEO-URL] durch die stabile, langfristig erreichbare
%   Adresse ersetzen (Uni-Cloud/Repository). Denselben Wert an BEIDEN
%   Stellen eintragen (Text-Link und \qrcode{...}).
% ============================================================

\chapter{Begleitmaterial}
\label{anh:begleitmaterial}

\section{Demonstrationsvideo}
\label{anh:video}

Als Begleitmaterial zu dieser Arbeit steht eine Videoaufzeichnung des Demonstrators zur Verfügung.
Sie zeigt eine frühe Ausbaustufe, in der die Szene über ein Menü in der virtuellen Umgebung manipuliert wird, das vordefinierte Befehle bereitstellt.
Das Video ist unter der folgenden Adresse abrufbar.

\begin{center}
\url{https://git.informatik.fh-nuernberg.de/jobboerse/forschung/bachelorarbeiten/bachelorarbeit-sandro-jaeger/-/blob/fe68f112f06df9da38eff687fd8df8460061186c/BA_VR_Demo_with_menu.mp4}

\vspace{1em}
\qrcode[height=3cm]{https://git.informatik.fh-nuernberg.de/jobboerse/forschung/bachelorarbeiten/bachelorarbeit-sandro-jaeger/-/blob/fe68f112f06df9da38eff687fd8df8460061186c/BA_VR_Demo_with_menu.mp4}

\vspace{0.5em}
{\footnotesize Abbildung: QR-Code zur Videoaufzeichnung.}
\end{center}